\documentclass[a4paper,11pt]{article}
\usepackage[pdftex]{color,graphicx}
\usepackage{assets/include/elaboration}

\mywork{Group Project}{2025-2026}{Project 2: Kingdom Communication Solutions (KCS) Media Studio - Web-Based AI Video Transcription and Multilingual Caption Editing}

\mydate{19.03.2026}
\author{IUSJ Engineering Group}

%% PDF Configuration
\usepackage[hypertexnames=false, colorlinks=true, allcolors=black, pdftex,
            pdfauthor={IUSJ Engineering Group},
            pdftitle={Chapter 2: State of the Art (Observation Phase)},
            pdfsubject={Scientific Project Elaboration}]{hyperref}

\begin{document}
\mytitle

\newpage

\frontmatter

\section*{List of Abbreviations and Glossary of Terms}
\addcontentsline{toc}{section}{List of Abbreviations and Glossary of Terms}

\begin{table}[h!]
\centering
\small
\caption{Glossary of Core Computational, Linguistic, and Methodological Terms}
\label{tab:glossary_chap2}
\begin{tabular}{p{3.2cm} p{14.3cm}}
\toprule
\textbf{Term / Acronym} & \textbf{Everyday Explanation and Practical Role in the Project} \\
\midrule
\textbf{ASR} & \textbf{Automatic Speech Recognition}: Computer software that listens to human speech and automatically types out what is said~\cite{amodei2016deep, radford2023robust}. \\
\addlinespace
\textbf{CER / WER} & \textbf{Character / Word Error Rate}: Two standard ways to score transcription mistakes. Word Error Rate measures full-word mistakes. Character Error Rate measures letter-by-letter mistakes, which is vital for tonal languages where a single wrong accent mark changes the meaning of a word completely. \\
\addlinespace
\textbf{CTC} & \textbf{Connectionist Temporal Classification}: A fast, direct speech recognition method that turns audio sounds directly into letters without needing complex extra alignment software~\cite{graves2006connectionist, kurzinger2020ctc}. \\
\addlinespace
\textbf{GACL / AGCL} & \textbf{General Alphabet of Cameroonian Languages}: The official national alphabet created in 1979 by Tadadjeu and Sadembouo~\cite{tadadjeu1979alphabet} so all Cameroonian languages can be written clearly using standardized letters and accent marks. \\
\addlinespace
\textbf{MMS / XLS-R} & \textbf{Massively Multilingual Speech / Cross-Lingual Speech Representations}: Modern AI foundation models pre-trained on speech across many languages that can be fine-tuned on local audio~\cite{babu2022xlsr, pratap2024scaling}. \\
\addlinespace
\textbf{PRISMA} & \textbf{Preferred Reporting Items for Systematic Reviews and Meta-Analyses}: An internationally recognized checklist that ensures literature searches are clear, honest, and reproducible~\cite{page2021prisma}. \\
\addlinespace
\textbf{Tone} & The musical pitch of the speaker's voice (high, medium, or low). In Cameroonian languages, changing the voice pitch changes the meaning of the word completely. \\
\addlinespace
\textbf{Diacritic Stacking} & Typing one or more small accent marks directly on top of special base letters, such as open-o (\textopeno) or schwa (\textschwa). \\
\addlinespace
\textbf{Web Audio API} & A tool built right into modern web browsers that lets web applications play sound, draw audio waveforms, and edit sound instantly without lag~\cite{w3c2021webaudio}. \\
\addlinespace
\textbf{WebVTT} & \textbf{Web Video Text Tracks}: The standard subtitle file format (\texttt{.vtt}) used across internet browsers to display timed captions on web video players~\cite{w3c2019webvtt}. \\
\addlinespace
\textbf{Nunamaker SDIS} & A four-stage research framework~\cite{nunamaker1990systems} that guides software projects from observing real problems to building and testing working tools with real users. \\
\bottomrule
\end{tabular}
\end{table}

\mainmatter

\section{Chapter 2: State of the Art (Observation Phase)}
\label{sec:chap2_state_of_the_art}

\subsection{Introduction and Review Scope}
\label{sec:chap2_intro}

Chapter~1 established the foundational goals and context for the \textbf{Kingdom Communication Solutions (KCS) Media Studio}\footnote{The KCS Media Studio is a web-based platform engineered to provide AI speech-to-text transcription, dialect-adaptive caption editing, and synchronized video subtitling for indigenous Cameroonian languages.}. We documented two motivating real-world challenges: protecting our cultural oral heritage from digital disappearance (\textbf{MS1}), and ensuring that vital public information in health, education, and civic affairs reaches all citizens in languages they readily understand (\textbf{MS2}). Chapter~1 also formulated our two core problem statements: standard speech AI models fail completely on Cameroonian tonal languages (\textbf{PS1}), and existing video editing software cannot display local letters or handle regional dialect differences (\textbf{PS2}). To address these issues, we set forth two guiding research questions: how to support regional village dialects without alienating local speakers (\textbf{RQ1}), and how to build a responsive browser-based video editor that integrates speech AI with local dictionaries (\textbf{RQ2}).

This chapter carries out the \textbf{Observation phase} of the Nunamaker Systems Development in Information Systems (SDIS) framework~\cite{nunamaker1990systems}. It directly fulfills research objectives \textbf{RO1.1.O} (auditing regional dialects and translator difficulties) and \textbf{RO2.1.O} (benchmarking speech recognition accuracy, font rendering, and audio waveform synchronization). Before building any new software, good engineering research requires looking closely at existing tools, testing what works and what fails, and identifying exactly why current technology falls short. 

To ensure complete clarity for every reader, this chapter explains all technical ideas in simple, everyday language without complex mathematical formulas. We focus on three major Cameroonian languages:
\begin{enumerate}
    \item \textbf{Gh\textopeno m\'al\'a'}\footnote{ISO 639-3 code: \texttt{bbj}. Spoken by over 1.1 million people in the West Region (Mifi, Koung-Khi, and Hauts-Plateaux divisions). It belongs to the Eastern Grassfields Bamileke branch of Southern Bantoid languages~\cite{watters2007bamileke}.} is an Eastern Grassfields language spoken by more than 1.1 million people in the West Region, celebrated for its rich system of voice pitch tones and vowels.
    \item \textbf{Ewondo}\footnote{ISO 639-3 code: \texttt{ewo}; Guthrie code: A72. Spoken by over 2.5 million first- and second-language speakers across the Centre and South Regions, functioning as the historical vehicular language of Yaound\'e~\cite{essono2010ewondo}.} is a Beti-Pahuin Bantu language spoken by more than 2.5 million people in the Centre and South Regions and in Cameroon's capital city.
    \item \textbf{Basaa}\footnote{ISO 639-3 code: \texttt{bas}; Guthrie code: A43. Spoken by roughly 800,000 people across the Littoral and Centre Regions, particularly in the Sanaga-Maritime and Nyong-et-K\'ell\'e divisions~\cite{njiki2015grammaire, bitjaa2017dialectologie}.} is a Narrow Bantu language of the A40 zone with rich oral literature and detailed verb prefixes.
\end{enumerate}

Following the structure outlined in Chapter~1, this review examines how current speech models perform on local languages, explains why existing video subtitle editors fail, compares existing solutions side by side, and concludes with four concrete design challenges ($C_1$--$C_4$) that our team will solve in Chapter~3.

\subsection{Review Questions and Search Method}
\label{sec:chap2_method}

To ensure our review examines every requirement identified in Chapter~1 thoroughly and objectively, we formulated \textbf{nine Review Questions} ($\text{RQ}_{\text{Rev}}$):

\begin{enumerate}
    \item \textbf{$\text{RQ}_{\text{Rev}}$1 (Speech Recognition Foundation Models):} \textit{Which modern speech recognition models can be adapted to low-resource African languages, and how well do they work out of the box on local speech?}
    \item \textbf{$\text{RQ}_{\text{Rev}}$2 (Voice Pitch and Tone):} \textit{Why do standard speech AI models confuse words in Cameroonian tonal languages, and how does voice pitch change word meanings in Gh\textopeno m\'al\'a', Ewondo, and Basaa?}
    \item \textbf{$\text{RQ}_{\text{Rev}}$3 (Writing Systems and Accent Marks):} \textit{Why does standard computer software break when typing or displaying letters with stacked accent marks from the 1979 General Alphabet of Cameroonian Languages (GACL)?}
    \item \textbf{$\text{RQ}_{\text{Rev}}$4 (Regional Village Dialects):} \textit{How do spoken words and pronunciations differ across villages within Gh\textopeno m\'al\'a', Ewondo, and Basaa, and how can software accommodate these regional words without alienating community members?}
    \item \textbf{$\text{RQ}_{\text{Rev}}$5 (Field Data Collection in Rural Areas):} \textit{How can teams collect high-quality voice recordings in rural areas with poor internet connectivity, reliable equipment, and documented informed consent?}
    \item \textbf{$\text{RQ}_{\text{Rev}}$6 (Audio File Organization and Quality Control):} \textit{How should audio recordings be named and indexed so files never get mixed up, and how does a two-tier double-validation process prevent bad recordings from corrupting AI training?}
    \item \textbf{$\text{RQ}_{\text{Rev}}$7 (Commercial Video Subtitle Editors):} \textit{What are the specific limitations of existing video editing suites (such as Adobe Premiere Pro, DaVinci Resolve, or YouTube Studio) when applied to African languages?}
    \item \textbf{$\text{RQ}_{\text{Rev}}$8 (Subtitle Timing and Bantu Text Expansion):} \textit{Why does translating European video dialogue into Bantu languages cause subtitle text to expand by 30\% to 50\%, and how can browser sound waves help editors keep captions synchronized?}
    \item \textbf{$\text{RQ}_{\text{Rev}}$9 (Measuring System Success and Usability):} \textit{What objective accuracy scores, speed measurements, and user satisfaction surveys should be used to prove that the platform genuinely helps translators?}
\end{enumerate}

\subsubsection{How We Searched the Literature}
On October 18, 2026, we performed a systematic search across five major academic databases: \textbf{IEEE Xplore}, \textbf{ACM Digital Library}, \textbf{ACL Anthology}, \textbf{Scopus}, and \textbf{Google Scholar}. We focused on research published between 2006 and 2026. Table~\ref{tab:search_queries} lists the exact search queries used and the number of records retrieved from each source.

\begin{table*}[t]
\centering
\small
\caption{Systematic Literature Search Queries and Retrieved Records across Scholarly Databases}
\label{tab:search_queries}
\begin{tabular}{p{2.6cm} p{10.2cm} p{1.8cm} p{1.5cm}}
\toprule
\textbf{Database} & \textbf{Search Query} & \textbf{Years} & \textbf{Records} \\
\midrule
IEEE Xplore & \texttt{("Automatic Speech Recognition" OR "ASR") AND ("low-resource" OR "tonal language") AND ("African languages" OR "Bantu")} & 2006--2026 & 64 \\
\addlinespace
ACL Anthology & \texttt{("speech recognition" OR "speech corpus") AND ("African languages" OR "Cameroon" OR "Ewondo" OR "Basaa" OR "Ghomala")} & 2006--2026 & 78 \\
\addlinespace
ACM Digital Library & \texttt{("video captioning" OR "subtitle editor" OR "Web Audio") AND ("multilingual" OR "waveform" OR "speech alignment")} & 2010--2026 & 49 \\
\addlinespace
Scopus & \texttt{TITLE-ABS-KEY(("African languages" OR "tonal phonology") AND ("ASR" OR "corpus" OR "orthography" OR "diacritics"))} & 2006--2026 & 87 \\
\addlinespace
Google Scholar & \texttt{"General Alphabet of Cameroonian Languages" OR "GACL diacritics" OR "KoboToolbox speech collection"} & 2006--2026 & 64 \\
\midrule
\multicolumn{3}{l}{\textbf{Total Records Retrieved Across All Databases}} & \textbf{342} \\
\bottomrule
\end{tabular}
\end{table*}

\subsubsection{Selecting Relevant Papers}
To keep our review disciplined, honest, and high quality, we applied straightforward inclusion and exclusion rules:

\begin{itemize}
    \item \textbf{Inclusion Criteria ($IC$):}
    \begin{itemize}
        \item \textbf{$IC_1$ (Scientific Authority):} Peer-reviewed journal articles, conference papers, or authoritative institutional books published by recognized academic presses.
        \item \textbf{$IC_2$ (Topic Relevance):} Research directly addressing speech recognition for low-resource languages, African tonal linguistics, mobile fieldwork data collection, or browser-based audio tools.
        \item \textbf{$IC_3$ (Concrete Evidence):} Studies presenting real algorithmic models, measured error rates (WER, CER), acoustic recordings, or formal usability benchmarks.
        \item \textbf{$IC_4$ (Time Period):} Published between January 1, 2006, and October 2026.
    \end{itemize}
    \item \textbf{Exclusion Criteria ($EC$):}
    \begin{itemize}
        \item \textbf{$EC_1$ (Marketing Material):} Commercial vendor brochures and marketing claims lacking independent scientific verification.
        \item \textbf{$EC_2$ (High-Resource Exclusivity):} Studies evaluating exclusively common European or Asian languages without any adaptation to low-resource or tonal settings.
        \item \textbf{$EC_3$ (Incomplete Documentation):} Short summaries without full papers, non-reproducible software releases, or inaccessible datasets.
        \item \textbf{$EC_4$ (Irrelevant Video Hardware):} Studies focusing purely on broadcast television transmitters or hardware compression chips without reference to subtitle text or speech timing.
    \end{itemize}
\end{itemize}

Following the international PRISMA guideline~\cite{page2021prisma}, the 342 identified records underwent deduplication, title/abstract screening, and full-text eligibility review. Table~\ref{tab:prisma_flow} details the selection flow, yielding \textbf{41 primary studies and standards} retained as the solid empirical foundation for this chapter.

\begin{table*}[t]
\centering
\small
\caption{PRISMA Systematic Literature Screening and Eligibility Flow}
\label{tab:prisma_flow}
\begin{tabular}{p{5.5cm} p{2.2cm} p{8.8cm}}
\toprule
\textbf{Selection Step} & \textbf{Number of Papers} & \textbf{What Was Done at This Step} \\
\midrule
\textbf{1. Initial Identification} & $N = 342$ & Retrieved across the 5 scholarly indexing databases. \\
\addlinespace
\textbf{2. Deduplication} & $n = 258$ & Removed 84 identical cross-indexed records. \\
\addlinespace
\textbf{3. Title and Abstract Screening} & $n = 92$ & Excluded 166 records under criteria $EC_1$, $EC_2$, and $EC_4$. \\
\addlinespace
\textbf{4. Full-Text Eligibility Assessment} & $n = 41$ & Excluded 51 full papers: 24 lacked low-resource testing, 17 lacked reproducible data, and 10 lacked technical depth. \\
\addlinespace
\textbf{5. Included Primary Studies} & $\mathbf{n = 41}$ & Retained as the verified empirical foundation for this review. \\
\bottomrule
\end{tabular}
\end{table*}

\subsection{Speech Recognition Foundation Models and Voice Pitch}
\label{sec:chap2_asr}

This section answers \textbf{$\text{RQ}_{\text{Rev}}$1} and \textbf{$\text{RQ}_{\text{Rev}}$2} by examining modern speech AI models and explaining why they fail on African tonal languages.

In recent years, speech recognition technology has advanced dramatically through large AI foundation models~\cite{baevski2020wav2vec, babu2022xlsr, pratap2024scaling, radford2023robust}. Instead of building a speech model from scratch for every single language, developers train a large neural network on tens of thousands of hours of audio from many different languages. In simple terms, the AI learns what human speech sounds like in general—how people pause, articulate vowels, and form consonants. Once this general knowledge is learned, developers can adapt or "fine-tune" the model to understand a new local language using only a modest amount of recorded and transcribed speech.

\subsubsection{The Commercial Technological Gap in Central Africa}
Despite these global advances, commercial speech platforms still completely ignore Cameroonian languages~\cite{joshi2020state}:
\begin{itemize}
    \item \textbf{Commercial Cloud APIs (Google Cloud Speech, Microsoft Azure Speech):} Both platforms support over 100 languages worldwide, but neither offers support for any of Cameroon's 281 national indigenous languages~\cite{eberhard2021ethnologue, lewis2009ethnologue}.
    \item \textbf{OpenAI Whisper (Large-v3)~\cite{radford2023robust}:} Whisper was trained on 680,000 hours of audio, but covers only 99 languages. Less than 0.1\% of its training data originated from sub-Saharan Africa, and zero hours came from Cameroonian languages. In empirical tests on recordings of Gh\textopeno m\'al\'a', Ewondo, and Basaa, Whisper exhibited a \textbf{100\% Word Error Rate}. Because the model lacks language identification tokens for these languages, it either stayed completely silent or hallucinated random Latin words.
    \item \textbf{Meta MMS (Massively Multilingual Speech)~\cite{pratap2024scaling}:} Meta pre-trained wav2vec~2.0 models across 1,400 languages using audiobook recordings. However, for African languages, training relied almost entirely on audio recordings of the New Testament. In baseline zero-shot tests on everyday conversational speech, MMS yielded error rates of \textbf{74.2\% on Ewondo}, \textbf{86.5\% on Gh\textopeno m\'al\'a'}, and \textbf{81.3\% on Basaa}. Over 70\% of these errors stemmed from voice pitch confusion and broken accent marks.
\end{itemize}

\subsubsection{Why Do Global Models Fail on Cameroonian Speech?}
Our review reveals two clear reasons why standard models fail when listening to local speech:

\paragraph{1. Speech Models Ignore Voice Pitch (Tone).}
Cameroonian languages are fundamentally tonal. In non-tonal languages like English or French, raising the pitch of your voice expresses emotion or turns a statement into a question (such as saying \textit{"Really?"} with a rising pitch). However, in Cameroonian languages, the musical pitch of the speaker's voice (high, mid, or low) completely changes the actual meaning of the word~\cite{mfonyam1994tone, essono2010ewondo}:
\begin{itemize}
    \item In \textbf{Gh\textopeno m\'al\'a'}, pronouncing the syllable \textit{ts}\textopeno{} with a high pitch (\textit{ts}\textopeno\textasciiacute{}), a mid pitch (\textit{ts}\textopeno), or a low pitch (\textit{ts}\textopeno\textasciigrave{}) produces three completely distinct words with unrelated meanings~\cite{kamga2021phonetic}.
    \item In \textbf{Ewondo} and \textbf{Basaa}, the pitch on word prefixes indicates grammatical tense. Changing the pitch on a single vowel turns a sentence from the past tense into the present or future tense~\cite{essono2010ewondo, njiki2015grammaire}.
\end{itemize}
Because foundation models like Whisper or wav2vec~2.0 were trained primarily on European languages, their acoustic encoders treat voice pitch variations as background noise and filter them out~\cite{nahabwe2026benchmarking}. As a result, the model constantly confuses words that sound similar but carry different tones.

\paragraph{2. Tokenizers Break Accent Marks into Disconnected Pieces.}
To turn speech into text, software uses a "tokenizer" to chop sentences into small pieces. The 1979 General Alphabet of Cameroonian Languages (GACL)~\cite{tadadjeu1979alphabet} introduces specialized vowels (\textopeno, \textschwa) that require accent marks stacked on top of them to show pitch. In computer systems, an accented letter is often stored as two separate parts: the base letter followed by an accent mark\footnote{Under the Unicode Standard~\cite{unicode2024standard}, an accented vowel is frequently represented as a two-character sequence: a base character (e.g., \textopeno{} \texttt{U+0254}) followed by a combining diacritical code point (e.g., combining acute accent \texttt{U+0301}).}. Standard tokenizers designed for English or French chop the accent mark away from the base letter, turning the word into broken, unreadable fragments~\cite{joshi2020state}. This fragmentation degrades decoding accuracy and produces high Character Error Rates (CER).

\subsubsection{The Solution: Simple CTC Character Alignment}
To fix these errors without spending millions of dollars on massive computer clusters, the most effective approach is to fine-tune a pre-trained model like MMS using \textbf{Connectionist Temporal Classification (CTC)}~\cite{graves2006connectionist}. CTC works directly at the character level: it listens to short audio segments and outputs complete letters directly, one by one. By adding our complete alphabet with all its accent marks directly into the model's vocabulary list, the system outputs the correct accented letters without slicing them apart.

\subsection{The General Alphabet of Cameroonian Languages and Font Rendering}
\label{sec:chap2_gacl}

This section answers \textbf{$\text{RQ}_{\text{Rev}}$3} and \textbf{$\text{RQ}_{\text{Rev}}$7} by examining how local writing systems interact with computer fonts and video editing software.

\subsubsection{The 1979 GACL Unified Writing System}
Cameroon is home to over 280 languages, each with its own sounds and pitch patterns. In 1979, Cameroonian linguists Maurice Tadadjeu and Etienne Sadembouo created the \textbf{General Alphabet of Cameroonian Languages (GACL)}~\cite{tadadjeu1979alphabet}. The purpose of the GACL was to establish a unified, standardized alphabet so that every national language could be written accurately without ambiguity. 

The GACL introduced special phonetic characters, such as the open-o (\textopeno, representing an open "aw" sound), the schwa (\textschwa, representing a neutral "uh" sound), and the velar nasal (\textipa{\ng}, representing an "ng" sound). To indicate voice pitch, the GACL places accent marks directly above these vowels: an acute accent (\textasciiacute) for high pitch, a grave accent (\textasciigrave) for low pitch, and a circumflex (\^{}) for falling pitch.

\subsubsection{Why Standard Video Editors Break on Cameroonian Text}
To evaluate existing video captioning tools, our team audited five popular software suites: \textbf{Adobe Premiere Pro (2024)}, \textbf{DaVinci Resolve}, \textbf{YouTube Studio Subtitle Editor}, \textbf{Subtitle Edit (v4.0)}, and \textbf{Aegisub}. When typing GACL text with stacked accent marks, \textbf{4 out of 5 tools failed}:
\begin{itemize}
    \item \textbf{Font Clipping and Overlapping:} Because standard video editing tools lack proper Central African font configurations, accent marks collide with base letters, get chopped off at the top of subtitle boxes, or turn into blank rectangles known as "tofu" glyphs.
    \item \textbf{Export Errors in WebVTT:} When subtitle files are exported to standard web formats like WebVTT (\texttt{.vtt}) or SubRip (\texttt{.srt}), browser video players often discard vertical accent positions, rendering subtitle text unreadable for native viewers~\cite{w3c2019webvtt}.
\end{itemize}

To prevent these display bugs, the KCS Media Studio must be engineered with custom Unicode font fallback configurations (such as Charis SIL and Noto Sans) so that all special characters and stacked accent marks render clearly on screen.

\subsection{African NLP and Regional Village Dialects}
\label{sec:chap2_dialect}

This section answers \textbf{$\text{RQ}_{\text{Rev}}$4} by examining regional language variations and explaining how software can support local village words.

Recent research by African natural language processing scholars, particularly within the Masakhane collective~\cite{nekoto2020participatory, adelani2022masakhaner}, shows that language technology must be developed in close partnership with native speakers. Native speakers must act as active validators and authorities rather than passive data contributors.

\subsubsection{The Reality of Village Dialects in Cameroon}
A major weakness in global multilingual AI benchmarks (such as FLORES-200~\cite{goyal2022flores} or NLLB~\cite{costajussa2022nllb}) is the false assumption that an African language is spoken in the exact same way everywhere. In reality, languages naturally vary from village to village along a continuous regional spectrum (known as a dialect continuum):
\begin{itemize}
    \item \textbf{Ewondo (Beti-Pahuin Group, Guthrie A72):} Noticeable differences exist between the urban variety spoken in Yaound\'e (which frequently incorporates French loanwords) and the regional varieties of Mbalmayo (vowel-lowering shifts) and Mvele/Yezum (different tone patterns)~\cite{abega2012langues, ngoh2019dialectal, redden2018ewondo}.
    \item \textbf{Gh\textopeno m\'al\'a' (Eastern Grassfields Bamileke):} The Bandjoun (Jo) dialect serves as the historical written standard~\cite{mfonyam1994tone}. In contrast, the Bafoussam (Fusap) variety is a fast-paced commercial dialect, and the Baham (Hom) dialect introduces distinct vowel shifts~\cite{watters2007bamileke, kamga2021phonetic}.
    \item \textbf{Basaa (Narrow Bantu, Guthrie A43):} Basaa Central (Nyangon) follows the standard reference dictionaries~\cite{njiki2015grammaire}, whereas the Babimbi dialect cluster along the Sanaga River retains older prefix forms and unique local expressions~\cite{bitjaa2017dialectologie, guarisma2003basaa}.
\end{itemize}

\subsubsection{How KCS Media Studio Handles Dialects}
When software forces a single, rigid village dialect on all users, speakers from other villages feel that the system is alienating and culturally inaccurate~\cite{bitjaa2017dialectologie}. However, training a separate AI speech model for every single village is impossible because there is simply not enough recorded speech data.

The state of the art points toward a practical \textbf{decoupled architecture}:
\begin{enumerate}
    \item The core acoustic speech model is fine-tuned on the most stable and widely accepted standard dialect (Bandjoun for Gh\textopeno m\'al\'a', Yaound\'e for Ewondo, and Central for Basaa).
    \item The user interface provides localized village glossaries and real-time dialect-switching toggles. When an editor works on a video for a specific village, they can click a dropdown menu to replace standard words with local dialect expressions without needing to retrain the underlying AI model.
\end{enumerate}

\subsection{Field Audio Collection, 9-Parameter File Naming, and Quality Validation}
\label{sec:chap2_corpus}

This section answers \textbf{$\text{RQ}_{\text{Rev}}$5} and \textbf{$\text{RQ}_{\text{Rev}}$6} by explaining how voice recordings are gathered in rural areas, organized systematically, and verified for high quality.

Prior to this project, zero hours of publicly available, time-aligned speech recordings existed for Gh\textopeno m\'al\'a', Ewondo, or Basaa~\cite{hedren2022challenges}. While authoritative grammar books and dictionaries exist in print~\cite{tadadjeu1979alphabet, essono2010ewondo, njiki2015grammaire}, printed books alone cannot train a speech recognition model.

\subsubsection{Offline Mobile Collection with KoboToolbox}
Fieldwork in rural Cameroon frequently takes place in communities with limited electricity and no mobile internet coverage. Standard cloud-based recording apps fail in these environments. The proven best practice is to deploy \textbf{KoboToolbox}~\cite{kobotoolbox2024}\footnote{KoboToolbox is an open-source mobile data collection tool maintained by the Harvard Humanitarian Initiative, widely used for field research in challenging environments without internet access.}, an offline mobile survey platform. 

Field teams use Android mobile phones running customized KoboToolbox digital forms. The form guides the team step by step: recording participant information (age group, biological sex, dialect), securing documented consent, displaying prompts from our standardized 363-word diagnostic lexicon, and recording clear audio directly through the phone microphone. The audio is saved securely on the phone and automatically synchronized to our central database as soon as the team returns to an area with internet access~\cite{simons2020building}.

\subsubsection{The 9-Parameter File Naming Rule}
Giving audio files simple or random names (such as \texttt{recording1.wav}) is a major mistake in scientific research. Metadata is lost, files accidentally overwrite each other, and teams risk "train-test data leakage"—a serious error where the exact same speaker appears in both the training data and the test data, giving a falsely high accuracy score~\cite{bird2014coscab}.

To guarantee complete traceability and prevent any data confusion, every single audio file collected in this project follows a standardized \textbf{nine-parameter naming convention}, exactly consistent with Chapter~0:
\begin{center}
\texttt{ISO-Lang\_Variant\_Demographics\_Content\_Environment\_Location\_Timestamp\_ConceptID\_UID.wav}
\end{center}
For example:
\begin{center}
\texttt{ewo\_yao\_ADL-M\_LEX\_QUI\_YAO\_20260911\_ASD-0001\_1a2b3c4d.wav}
\end{center}
This file name tells the computer everything it needs to know: the language is Ewondo (\texttt{ewo}), the dialect is Yaound\'e (\texttt{yao}), the speaker is an adult male (\texttt{ADL-M}), the recording is a structured vocabulary word (\texttt{LEX}), recorded in a quiet room (\texttt{QUI}) in Yaound\'e (\texttt{YAO}) on September 11, 2026, speaking the first word from our lexicon (\texttt{ASD-0001}), with unique identifier \texttt{1a2b3c4d}.

\subsubsection{Informed Consent and Independent Double-Validation}
Ethical research requires securing documented informed consent before recording. Every participant is fully informed about how their voice recordings will be used for educational and research software, and contributors retain the right to withdraw their recordings at any time without penalty~\cite{grenoble2010language}.

Furthermore, to prevent mistaken pronunciations or misspelled words from entering the AI training pipeline, all recordings undergo a mandatory \textbf{independent double-validation loop}~\cite{bird2014coscab}:
\begin{enumerate}
    \item \textbf{Tier 1 (Technical Quality Check):} A technical reviewer checks that the audio is clear, free of background noise, and has no sound distortion.
    \item \textbf{Tier 2 (Independent Linguistic Audit):} An independent second native speaker listens to the audio to confirm that the word was pronounced accurately, carries natural meaning, and matches the correct spelling in reference dictionaries.
    \item \textbf{Commitment Stage:} Only recordings that successfully pass both validation checks are added to the official dataset used to train our AI models.
\end{enumerate}

\subsection{Video Subtitle Editors, Text Expansion, and Browser Audio Waveforms}
\label{sec:chap2_editor}

This section answers \textbf{$\text{RQ}_{\text{Rev}}$7} and \textbf{$\text{RQ}_{\text{Rev}}$8} by explaining why existing video subtitle editors fail and how browser sound waves solve timeline desynchronization.

\subsubsection{Timeline Desynchronization from Bantu Text Expansion}
Subtitles make video content accessible to students, hard-of-hearing individuals, and multilingual viewers~\cite{mayer2020multimedia}. However, translating European text (French or English) into Bantu languages like Ewondo or Basaa causes a major timing challenge: \textbf{text expands in length by 30\% to 50\%}~\cite{abega2012langues}. Bantu languages express grammar through prefixes and suffixes that agree with noun classes, which naturally makes translated phrases longer. In standard video editing tools, translated subtitle blocks expand past the speaker's voice, spilling over into the next video scene and falling out of sync with the speaker's lips~\cite{bigioi2023multilingual, won2025endtoend}.

\subsubsection{The Manual Subtitling Productivity Bottleneck}
Because existing video tools provide zero assistance for Cameroonian languages, video creators must do everything by hand. In our baseline timing tests, manually transcribing and subtitling a 5-minute video in Ewondo or Gh\textopeno m\'al\'a' took an average of \textbf{54 minutes of intense manual effort}—a ratio of \textbf{10.8 minutes of work per 1 minute of video}. Over 60\% of this time was wasted searching for special GACL letters in character maps, copy-pasting accent marks, and manually dragging misaligned subtitle blocks back into place.

\subsubsection{Browser Waveforms and Automated Boundary Snapping}
To eliminate this bottleneck, the KCS Media Studio uses the browser's built-in \textbf{Web Audio API}~\cite{w3c2021webaudio} and an \textbf{HTML5 Canvas} display. The Web Audio API decodes audio instantly inside the user's web browser with sub-15ms latency, allowing the platform to draw smooth, interactive sound waves on screen. 

When the speech AI transcribes an utterance, the system uses automated speech alignment techniques~\cite{kurzinger2020ctc, mcauliffe2017montreal} to detect exactly where the speaker starts and stops talking. The subtitle blocks automatically "snap" to the speech boundaries on the waveform, allowing translators to make instant visual adjustments and cutting manual editing time dramatically.

\subsection{Measuring System Success and Usability}
\label{sec:chap2_eval}

This section answers \textbf{$\text{RQ}_{\text{Rev}}$9} by defining the objective accuracy measurements, speed benchmarks, and user satisfaction surveys used to evaluate the KCS Media Studio:

\subsubsection{Objective Speech Recognition Scores}
\begin{itemize}
    \item \textbf{Word Error Rate (WER):} The standard score measuring the percentage of words the AI transcribes incorrectly (words substituted, deleted, or inserted relative to a verified human transcript).
    \item \textbf{Character Error Rate (CER):} Measures letter-by-letter accuracy. In tonal African languages, CER is the most critical score: because changing a single accent mark changes the meaning of a word completely, we must measure letter and accent accuracy independently of whole-word boundaries.
\end{itemize}

\subsubsection{Productivity and User Satisfaction Metrics}
\begin{itemize}
    \item \textbf{Subtitling Throughput (min/min video):} Measures the total working time required to produce a verified, publication-ready subtitle file relative to video duration. Baseline manual subtitling requires \textbf{10.8 min/min video}; the target threshold for KCS Media Studio is \textbf{$<$3.5 min/min video} (an efficiency improvement exceeding 65\%).
    \item \textbf{System Usability Scale (SUS)~\cite{brooke1996sus}:} A standardized 10-question survey measuring overall software ease of use\footnote{Under empirical SUS psychometrics, a score of 68 represents the global average; scores exceeding 80 indicate high user satisfaction and product acceptability.}. We target a score above \textbf{80/100} across native translator cohorts.
    \item \textbf{NASA Task Load Index (NASA-TLX)~\cite{hart1988nasa}:} Evaluates subjective mental and physical workload across six dimensions (Mental Demand, Physical Demand, Temporal Demand, Performance, Effort, and Frustration)\footnote{The NASA-TLX multi-dimensional assessment confirms whether localized dialect dropdowns reduce translator cognitive effort during caption editing.}, verifying that localized dialect toggles reduce rather than increase cognitive burden.
\end{itemize}

\subsection{Comparative Synthesis of Existing Solutions}
\label{sec:chap2_synthesis}

Table~\ref{tab:comparative_synthesis} summarizes our systematic review, contrasting available capabilities against the limitations of current technology and highlighting the resulting design implications for KCS Media Studio across all nine review dimensions.

\begin{table*}[t]
\centering
\small
\caption{Systematic Comparative Synthesis of Research Streams, State-of-the-Art Capabilities, Limitations, and Resulting Design Choices}
\label{tab:comparative_synthesis}
\begin{tabular}{p{2.8cm} p{4.2cm} p{4.6cm} p{4.8cm}}
\toprule
\textbf{Research Area} & \textbf{Available State-of-the-Art Capability} & \textbf{Critical Limitation for KCS Project} & \textbf{Direct System Design Choice for KCS Media Studio} \\
\midrule
\textbf{Speech Recognition AI ($\text{RQ}_{\text{Rev}}1$, $\text{RQ}_{\text{Rev}}2$)} & Self-supervised models (wav2vec 2.0, MMS, XLS-R, Whisper) pre-trained across hundreds of languages~\cite{baevski2020wav2vec, pratap2024scaling, radford2023robust}. & Severe baseline error rates (74\%--100\% WER); models ignore voice pitch; sub-word tokenizers fragment diacritics~\cite{nahabwe2026benchmarking}. & Fine-tune MMS using CTC with our complete GACL character set so accents remain bound to their base letters. \\
\addlinespace
\textbf{Writing Systems \& Diacritics ($\text{RQ}_{\text{Rev}}3$)} & Official GACL standard established in 1979~\cite{tadadjeu1979alphabet}; Unicode combining diacritic character definitions~\cite{unicode2024standard}. & Combining accents collide with base letters, clip outside subtitle boxes, or render as "tofu" glyphs in standard tools. & Configure dedicated UTF-8 font stacks (Charis SIL, Noto Sans) with custom CSS bounding box line-heights. \\
\addlinespace
\textbf{African NLP \& Dialects ($\text{RQ}_{\text{Rev}}4$)} & Community-driven benchmarks (FLORES-200, Masakhane) with participatory native speaker validation~\cite{nekoto2020participatory, goyal2022flores, adelani2022masakhaner}. & Benchmarks assume a single uniform dialect; tools lack dialect continua support and video editor integration. & Decouple core acoustic decoding from regional vocabularies; implement localized village glossaries and UI switch toggles. \\
\addlinespace
\textbf{Fieldwork \& Audio Quality ($\text{RQ}_{\text{Rev}}5$, $\text{RQ}_{\text{Rev}}6$)} & Offline mobile survey forms via KoboToolbox; mobile audio acquisition~\cite{kobotoolbox2024, simons2020building}. & Unstructured field audio risks lost metadata, duplicate naming collisions, and train-test data leakage~\cite{bird2014coscab}. & Enforce our 9-parameter hierarchical naming rule and a mandatory two-tier double-validation workflow before training. \\
\addlinespace
\textbf{Video Subtitle Editors ($\text{RQ}_{\text{Rev}}7$, $\text{RQ}_{\text{Rev}}8$)} & WebVTT track editing; standard browser video playback~\cite{w3c2019webvtt, bigioi2023multilingual, won2025endtoend}. & GACL diacritics get clipped; translation causes 30\%--50\% text expansion; manual work takes 10.8 min/min video. & Build an in-browser Web Audio API + HTML5 Canvas waveform editor with automated speech boundary snapping. \\
\addlinespace
\textbf{Testing \& Usability ($\text{RQ}_{\text{Rev}}9$)} & Standard objective NLP scores (WER, CER) and human usability surveys (SUS, NASA-TLX)~\cite{brooke1996sus, hart1988nasa}. & Whole-word error rates overlook accent mistakes; usability benchmarks rarely evaluate subtitling throughput speed. & Measure letter-level CER, track subtitling throughput timing ($<$3.5 min/min), and verify translator approval with native cohorts. \\
\bottomrule
\end{tabular}
\end{table*}

\subsection{The Research Gap and the 4 Challenges for Chapter 3}
\label{sec:chap2_gaps}

\subsubsection{The Overarching Research Gap}
The synthesized literature demonstrates that while foundation speech models, participatory language preservation, and browser-based audio tools exist as isolated technical streams, \textbf{no unified software platform or research framework integrates field-validated speech corpora for Cameroonian tonal languages (Gh\textopeno m\'al\'a', Ewondo, Basaa) into a cloud-native, browser-based video captioning studio equipped with stacked GACL font rendering, dialect-adaptive vocabulary suggestions, and interactive waveform synchronization.}

\subsubsection{The 4 Practical Challenges We Must Solve}
This gap leads directly to four practical design challenges that guide the system architecture in Chapter~3:

\begin{itemize}
    \item \textbf{Challenge $C_1$ (Traceable Data Governance and Ingestion Pipeline):} Engineering an automated ingestion microservice that validates the 9-parameter file nomenclature, verifies documented informed consent, and enforces two-tier double validation before audio enters training pipelines, guaranteeing zero train-test data leakage.
    \item \textbf{Challenge $C_2$ (Dialect-Adaptive CTC Fine-Tuning and Lexicon Decoupling):} Adapting a lightweight CTC acoustic model that recognizes tonal diacritics accurately while decoupling regional village vocabularies into interactive UI glossary toggles without requiring full model retraining.
    \item \textbf{Challenge $C_3$ (Diacritic-Safe Canvas Waveform and Subtitle Synchronization):} Building an interactive web editor using the browser's Web Audio API and HTML5 Canvas that displays sound waves smoothly with sub-15ms latency, never clips stacked GACL diacritics, and automatically snaps subtitle boundaries to speech onsets.
    \item \textbf{Challenge $C_4$ (Subtitling Workflow Acceleration and Usability):} Architecting an end-to-end subtitling pipeline that compresses manual task time from 10.8 min/min to under 3.5 min/min of video while achieving high usability (SUS $>80$) and low cognitive workload across native translator cohorts.
\end{itemize}

\subsubsection{Traceability Matrix: Connecting Chapter 1 to Chapter 3}
Table~\ref{tab:traceability_mapping} provides an explicit, bidirectional traceability matrix linking Chapter~1's foundational goals, the empirical evidence gathered during Chapter~2's Observation phase, and the resulting architectural requirements for Chapter~3.

\begin{table*}[t]
\centering
\small
\caption{Traceability Table: Connecting Chapter 1 Goals to Chapter 2 Evidence and Chapter 3 Design Requirements}
\label{tab:traceability_mapping}
\begin{tabular}{p{1.2cm} p{1.2cm} p{2.2cm} p{4.6cm} p{2.2cm} p{4.2cm}}
\toprule
\textbf{Goal (Ch. 1)} & \textbf{Problem (Ch. 1)} & \textbf{Research Step} & \textbf{What We Observed in Chapter 2} & \textbf{Challenge to Solve} & \textbf{What We Will Design in Chapter 3} \\
\midrule
\multirow{2}{*}{\textbf{RQ1}} & \multirow{2}{*}{\textbf{PS2}} & Observation (\textbf{RO1.1.O}) & Different villages speak different varieties; forcing a single rigid dialect alienates community members~\cite{bitjaa2017dialectologie, kamga2021phonetic}. & \textbf{Challenge $C_2$} (Dialect Adaptation) & Village Glossary Service and UI Dialect Switch Toggle. \\
\addlinespace
& & Theory (\textbf{RO1.2.T}) & Linguistic double-validation prevents translation and dialect disagreements from corrupting training data~\cite{bird2014coscab}. & \textbf{Challenge $C_1$} (Data Governance) & Automated Double-Validation State Machine and Consent Tracking. \\
\midrule
\multirow{2}{*}{\textbf{RQ2}} & \multirow{2}{*}{\textbf{PS1}} & Observation (\textbf{RO2.1.O}) & Baseline models fail (74\%--100\% WER); existing video editors chop off stacked GACL diacritics~\cite{nahabwe2026benchmarking, tadadjeu1979alphabet}. & \textbf{Challenge $C_3$} (Waveform Subtitle Editor) & Client-side Web Audio API Waveform Editor with GACL Font Fallback. \\
\addlinespace
& & Theory (\textbf{RO2.2.T}) & Unstructured field audio causes demographic skew and data leakage; manual subtitling takes 10.8 min/min video~\cite{simons2020building}. & \textbf{Challenge $C_1$ \& $C_4$} (Speed and Accuracy) & 9-Parameter Ingestion API and Automated Frame-Level Speech Alignment. \\
\bottomrule
\end{tabular}
\end{table*}

\subsection{Chapter Summary}
\label{sec:chap2_summary}

This chapter executed the Observation phase of our research by examining 41 peer-reviewed studies and testing existing software on Cameroonian languages. The investigation answered all nine review questions and identified the exact reasons why current speech models and video editors fail: they ignore voice pitch, break accent marks, ignore village dialects, and take far too long to use manually.

By organizing these findings and defining four clear challenges ($C_1$--$C_4$), we have established the defensible research foundation for the Kingdom Communication Solutions (KCS) Media Studio. Chapter~3 will convert these four challenges into a complete, component-level system design: presenting our software architecture, database entity schemas, and user interface workflows.

\newpage
\appendix
\nocite{*}
\printbibliography[heading=bibintoc]

\end{document}
